% --- SEZIONE 2: MODELLO DI MEMORIA E STACK ---
\section{Ottimizzazione e Strutture Dati}

\subsubsection{Switch Case}
\begin{itemize}
    \item \textbf{Pro:} Anzahl der Taktzyklen definiert und vorhersehbar.
    \item \textbf{Contra:} Die Bedingung muss notwendigerweise eine Konstante vom Typ \texttt{int} sein.
    \item \textbf{Hinweis:} Die Reihenfolge der \texttt{case} ist für die Effizienz der Suche relevant.
\end{itemize}

\subsubsection{Flusso di Controllo in Assembly}
\begin{itemize}
    \item \textbf{Selektionen (If/Else):} Es wird ein \textbf{Sprung nach vorn} mit der \textbf{invertierten Bedingung} gegenüber C verwendet (wenn die Bedingung falsch ist, springt man über den \textit{then}-Block hinaus).
    \item \textbf{Iterationen (Loop):} Es wird ein \textbf{Sprung nach hinten} am Ende der Schleife verwendet, wobei die \textbf{Assembly-Bedingung wie in C} beibehalten wird (Sprung zum Anfang, wenn die Bedingung noch wahr ist).
\end{itemize}

% \subsubsection{Tecniche di Ottimizzazione}
% \begin{itemize}
%     \item \textbf{Loop Unrolling:} Tecnica che consiste nel ripetere il corpo del loop 
%     per ridurre l'overhead dei \textit{branch} e degli incrementi del contatore.
%     \texttt{\#pragma unroll (n)}
%     \item \textbf{Pointers (Stackless):} L'uso dei registri invece dello stack per manipolare i puntatori è significativamente più veloce (spesso più del doppio).
% \end{itemize}

\subsubsection{Efficienza delle Strutture Dati}
\begin{itemize}
    \item \textbf{Lokale Arrays:} Sehr \textbf{ineffizient}. 
    hohe Anzahl von Zyklen nur für die Initialisierung im Stack-Frame 
    \item \textbf{Strukturen:} Es ist empfehlenswert, die \textbf{Strukturen global zu definieren}. Die lokale Initialisierung auf dem Stack führt zu einem hohen Overhead an Taktzyklen.
    \item \textbf{Bitfields:} Extrem \textbf{ineffizient}. 
    Obwohl sie Speicher sparen, erfordern sie komplexe Operationen von
     \textit{Read-Modify-Write} (LDR, Shift, Mask, OR, STR)
\end{itemize}